\documentclass{UAmathtalk}

\author{Kristine Bauer}
\address{University of Calgary, Canada}
\urladdr{https://profiles.ucalgary.ca/kristine-bauer}
\title{The Categorical Derivative in~Functor~Calculus}
\date{Thursday, September 17, 2026}

\renewcommand*{\where}{Massry B008}
\renewcommand*{\abstractsize}{\normalsize}


\begin{document}

\maketitle

\begin{abstract}
In algebraic topology and elsewhere in mathematics, functors serve as invariants: they let us tell objects apart even when those objects are linked by weak equivalences, which can otherwise obscure genuine differences. Such functorial invariants are frequently very hard to compute. Functor calculus, introduced by Goodwillie in the 1990s, gives a systematic way to attack this problem: to a functor it assigns a sequence of ``degree~$n$'' approximations, in close analogy with the Taylor polynomials of a function of a real variable. 

Given how far the analogy with Taylor series reaches, it is natural to ask whether functor calculus is genuinely an instance of some general notion of calculus. In this talk I will describe one precise answer. Category theorists --- Blute, Cockett, Cruttwell, Rosický, Seely, and others --- have developed an axiomatic theory of differentiation that generalizes the derivative of a smooth map between manifolds, organized around the notions of cartesian differential category and tangent category. Bauer, Johnson, Osborne, Riehl, and Tebbe showed that abelian functor calculus (one of several varieties of the theory) supplies the derivative making the category of abelian categories a cartesian differential category, up to homotopy. Building on this, Bauer, Burke, and Ching showed that the Goodwillie derivative provides exactly the structure needed to exhibit the $\infty$-category of $\infty$-categories as a tangent $\infty$-category.

I will give enough background on functor calculus and on these categorical notions, with examples throughout, to keep the talk accessible to a general mathematical audience.
\end{abstract}

\end{document}
